% with-edges.tex — fixture: paper that uses every edge type.
%
% Purpose: exercise the v0.2 pipe-format edge markers (RRP-0002).

\documentclass{rrxiv}

\rrxivid{rrxiv-fixture-with-edges}
\rrxivversion{v1}
\rrxivprotocolversion{0.1.0}
\rrxivlicense{CC-BY-4.0}
\rrxivtopics{example,conformance}

\title{Conformance fixture: paper with all four edge kinds}
\author{rrxiv Project}

\begin{document}
\maketitle

\begin{abstract}
Three claims plus four inline edges that cover every edge kind:
\texttt{depends\_on}, \texttt{supports}, \texttt{extends}, and
\texttt{contradicts}.
\end{abstract}

\section{Foundation}

\begin{claim}[Foundation claim]
\label{claim:foundation}
A foundational claim other claims will reference.
\end{claim}

\section{Derived claims}

\begin{claim}[Depends on the foundation]
\label{claim:derived}
A claim that builds on the foundation. The dependency is declared
explicitly via \texttt{\textbackslash dependson}.
\end{claim}

\dependson{rrxiv-fixture-with-edges:claim:derived}{rrxiv-fixture-with-edges:claim:foundation}

\begin{claim}[Supports the foundation]
\label{claim:replication}
A claim that supports the foundation as a kind of replication.
\end{claim}

\supports{rrxiv-fixture-with-edges:claim:replication}{rrxiv-fixture-with-edges:claim:foundation}

\section{Cross-paper edges}

\extendsclaim{rrxiv-fixture-with-edges:claim:foundation}{rrxiv-0001:claim:queryability}

\contradicts{rrxiv-fixture-with-edges:claim:derived}{some-other-paper:claim:opposite}

\end{document}
